% ============================================================================
\chapter{Бесфононный кремниевый гетеропереход \texorpdfstring{$120 + 8 = 128$}{120 + 8 = 128}, общая конструкция предела по счётным моноидам стадий и гёделевская неизбежность лабораторного арбитража}
\label{ch:mono-solar-photovoltaics}
\label{sec:mono-epistemological-closure}
% ============================================================================

\section{Бесфононный кремниевый гетеропереход \texorpdfstring{$120 + 8 = 128$}{120 + 8 = 128}, барьер пассивации Паули \texorpdfstring{$45^2 < 2\,048 < 46^2$}{45^2 < 2048 < 46^2} и вычисления без тепла Ландауэра}

\begin{theorem}[Резонансный замок седла $120 + 8 = 128 = Q_\square(3)$, барьер пассивации Паули $45^2 < 2\,048 < 46^2$, окно Френеля $288/289$ и безветвевая кремниевая архитектура в коридоре Гейзенберга]
\label{thm:mono-silicon-photocascade-128}
В естественных безразмерных координатах носителя на четырёхслойной гетероструктуре «Кварц ($\mathrm{SiO_2}$) -- Тетрацен ($\mathrm{C_{18}H_{12}}$) -- Пассивирующий мост ($\mathrm{SiO_x}/\mathrm{Al_2O_3}$) -- Кремний ($\mathrm{IBC\text{-}Si}$, $Z(\mathrm{Si})=14$)» выполняются следующие точные квантово-инженерные законы (\texttt{complete\_\allowbreak silicon\_\allowbreak hetero\_\allowbreak engineering\_\allowbreak operator}, \texttt{singlet\_\allowbreak fission\_\allowbreak tetracene\_\allowbreak silicon\_\allowbreak saddle\_\allowbreak 128\_\allowbreak resonance}, \texttt{reversible\_\allowbreak computing\_\allowbreak zero\_\allowbreak landauer\_\allowbreak heat\_\allowbreak and\_\allowbreak paradox\_\allowbreak resolutions\_\allowbreak operator}):
\begin{enumerate}[label=\textbf{(\arabic*)},leftmargin=*]
\item \textbf{Резонансный замок седла четверти оборота $120 + 8 = 128 = Q_\square(3)$, барьер пассивации тетраэдра Паули $45^2 < 2\,048 < 46^2$ и антиотражающее окно Френеля $288/289$:}
  молекула тетрацена $\mathrm{C_{18}H_{12}}$ содержит $Z(\mathrm{C_{18}H_{12}}) = 18\times 6 + 12\times 1 = 120 = Q_\square(3) - |{\triangle}(3)|$ электронов и в сопряжении с координатным кислородным мостом $Z(\mathrm{O}) = |{\triangle}(3)| = 8$ попадает точно в седло четверти оборота:
  \begin{align}
    Z(\mathrm{C_{18}H_{12}}) + Z(\mathrm{O}) &= 120 + 8 \;=\; 128 \;=\; Q_\square(3), \label{eq:mono-tetracene-128-lock} \\
    \bigl(\mathcal{C}_3(128),\ \mathcal{S}_3(128)\bigr) &= (0,\ +128). \label{eq:mono-tetracene-128-cs}
  \end{align}
  При равномерном расщеплении максимального полупериодного импульса Паули $\Phi_{\mathrm{Pauli}}(3) = 256$ по $|O_2(3)| = 12$ лучам звезды кубооктаэдра фазовый шаг на один луч равен $\Delta\varphi_{\mathrm{ray}}(3) = \lfloor 256 / 12 \rfloor = 21 = \Delta\varphi_{\mathrm{seam}}(3) < 41 = \Delta\varphi_{\mathrm{disc}}(3)$, поэтому число излучаемых тепловых квантов на каждом луче тождественно равно нулю: $\hat{n}_E(21) = \lfloor 21 / 41 \rfloor = 0$. На интерфейсе оксидной пассивации непассивированный алюмосиликатный тетраэдр $Z(\mathrm{AlO_4}) = N_\varphi(3) + d_{\mathrm{in}}(\square)Z(\mathrm{O}) = 13 + 32 = 45$ лежит строго ниже полупорога плотности состояний Паули $P_{\mathrm{Pauli}}(3) := K_P(3)^2 / (c_0/c)^2 = 64^2 / 2 = 2\,048$, тогда как тетраэдр кремнезёма $Z(\mathrm{SiO_4}) = 14 + 32 = 46 = S_3(3)/2$ и водородно-пассивированный тетраэдр $Z(\mathrm{AlO_4H}) = 45 + 1 = 46$ запирают барьер Паули при безразмерном коэффициенте прохождения согласованного окна Френеля $T_{\mathrm{Fresnel}}(3) = 288/289$:
  \begin{equation}
    \label{eq:mono-882-masons-stothers}
    Z(\mathrm{AlO_4})^2 = 2\,025 \;<\; P_{\mathrm{Pauli}}(3) = 2\,048 \;<\; 2\,116 = Z(\mathrm{SiO_4})^2,
    \qquad
    T_{\mathrm{Fresnel}}(3) = \frac{|{\triangle}(3)|\,d_{\mathrm{in}}(\square)\,\mathrm{mod\_bits}(3)}{\bigl(\mathrm{mod\_bits}(3) + |{\triangle}(3)|\bigr)^2} = \frac{288}{289}.
  \end{equation}
\item \textbf{A/B-критерий порога Паули ($Z = 45$ против $Z = 46$) и безветвевой кремниевый вычислительный стек (\texttt{Ring9} / \texttt{Bus36} / \texttt{VerletCell72}):}
  в каталитическом и твердотельном A/B-тесте родий $Z(\mathrm{Rh}) = 45 = \mathrm{mod\_bits}(3)\,d_{\mathrm{out}}(E) < 46 = S_3(3)/2$ лежит строго ниже полуоболочки третьей сферы кубооктаэдра и сохраняет открытый канал переноса, тогда как палладий $Z(\mathrm{Pd}) = 46 = S_3(3)/2$ замыкает полуоболочку Паули. Аппаратный стек на регистрах \texttt{Ring9} ($9\text{ бит}$), \texttt{Bus36} ($B_1(3) = 36\text{ бит}$) и \texttt{VerletCell72} ($B_2(3) = 72\text{ бита}$) выполняет биективный шаг Верле~\eqref{eq:mono-verlet-law} без условных ветвлений (\texttt{0 branch mispredictions}) и при работе в коридоре $|\Delta\varphi| \le 40 < 41$ выделяет $\hat{n}_E = 0$ квантов тепла Ландауэра.
\end{enumerate}
\end{theorem}

\begin{proof}
В пункте~(1) $18\times 6 + 12 = 120 = 128 - 8$, $120 + 8 = 128 = Q_\square(3)$, $\mathcal{C}_3(128) = 128 - 128 = 0$, $\mathcal{S}_3(128) = +128$; на каждом из $12$ лучей $\lfloor 256/12 \rfloor = 21 < 41 \implies \hat{n}_E(21) = \lfloor 21/41 \rfloor = 0$; для тетраэдров пассивации $45^2 = 2\,025 < 64^2/2 = 2\,048 < 2\,116 = 46^2$, а для окна Френеля $8\times 4\times 9 = 288 = (9+8)^2 - 1 = 289 - 1$. В пункте~(2) $S_3(3) = 10\cdot 3^2 + 2 = 92$, откуда $S_3(3)/2 = 46 = Z(\mathrm{Pd}) > 45 = Z(\mathrm{Rh})$, а биективность шага Верле на $\mathbb{Z}_{512}^2$ доказана в Теореме~\ref{thm:mono-verlet-uniqueness}.
\end{proof}

\section{Общий оператор предела по счётным моноидам стадий, проективные башни Мальцева и двусторонняя эквивалентность носителей \texorpdfstring{$\mathcal{M} \leftrightarrow \mathcal{T}$}{M <-> T}}

В классической теории моделей А.\,И.~Мальцева (\emph{«Алгебраические системы»}, 1970) алгебраическая система $\mathfrak{A} = \langle A;\,\Omega_F,\Omega_P\rangle$ определяется своим \emph{носителем} (основным множеством сорта~$A$) и сигнатурой операций и предикатов. В построенной в Главах~1--7 двусортной системе $\langle\mathcal{M},\mathcal{T};\,\pi_{M\to T},\,\iota_{T\to M}\rangle$ общая конструкция предела (теорема \texttt{thm17\_\allowbreak model\_\allowbreak theoretic\_\allowbreak termination\_\allowbreak kernel\_\allowbreak annihilation\_\allowbreak and\_\allowbreak godel\_\allowbreak tarski\_\allowbreak oracle\_\allowbreak necessity} в модуле \texttt{ModelTheoreticLimit.lean}) объединяет пять алгебраических механизмов без привлечения неконструктивной аксиомы выбора (\texttt{0 Classical.choice}):
\begin{enumerate}[label=\textbf{\arabic*.},leftmargin=*]
\item \textbf{Конгруэнция ядра гомоморфизма алгебраических систем Мальцева, элементарная эквивалентность $\mathrm{Th}(M_{\mathrm{ext}}) = \mathrm{Th}(M_{\mathrm{base}})$ и стабилизация проективной башни $\{M_n\}_{n\in\mathbb{N}}$ (\texttt{MalcevSystem}, \texttt{MalcevHom}, \texttt{malcevKernel}, \texttt{pullBackSatisfaction}, \texttt{towerProjectToBase}, \texttt{stageFold}):}
  для любого гомоморфизма алгебраических систем Мальцева $\pi : \langle A;\,\star_A,\mathrm{rel}_A\rangle \to \langle B;\,\star_B,\mathrm{rel}_B\rangle$, сохраняющего операции и предикаты, отношение эквивалентности по ядру $(x \equiv_{\ker\pi} x') \iff (\pi(x) = \pi(x'))$ является конгруэнцией по операциям и предикатам:
  \begin{equation}
    \label{eq:mono-malcev-congruence}
    x \equiv_{\ker\pi} x' \;\land\; y \equiv_{\ker\pi} y' \quad\Longrightarrow\quad (x \star_A y) \equiv_{\ker\pi} (x' \star_A y') \;\land\; \bigl(\mathrm{rel}_A(x) \leftrightarrow \mathrm{rel}_A(x')\bigr),
  \end{equation}
  а для отношения выполнимости формул Тарского--Мальцева $(M, x) \models \varphi$ вдоль любого расщепляющегося уширения $\pi_{\mathrm{ext}} : M_{\mathrm{ext}} \twoheadrightarrow M_{\mathrm{base}}$ ($\pi_{\mathrm{ext}}\circ\mathrm{sect} = \mathrm{id}_{M_{\mathrm{base}}}$) классы ядра элементарно неразличимы, и теории всех универсально истинных предложений на $M_{\mathrm{ext}}$ и на $M_{\mathrm{base}}$ совпадают тождественно:
  \begin{equation}
    \label{eq:mono-elementary-equivalence}
    \forall\varphi \in \mathrm{Formula}, \qquad \bigl(\forall y \in M_{\mathrm{ext}},\ (M_{\mathrm{ext}}, y) \models \varphi\bigr) \;\Longleftrightarrow\; \bigl(\forall m \in M_{\mathrm{base}},\ (M_{\mathrm{base}}, m) \models \varphi\bigr).
  \end{equation}
  Индуктивная свёртка любого конечного слова стадий $w = [s_1, \dots, s_r]$ вдоль действия $\mathrm{step} : \Sigma \to \mathcal{M} \to \mathcal{M}$, сохраняющего проекцию ($\forall s\in\Sigma,\ \pi(\mathrm{step}(s, x)) = \pi(x)$), удовлетворяет закону моноида $\mathrm{stageFold}(w_1\mathbin{+\!+}w_2, x) = \mathrm{stageFold}(w_2, \mathrm{stageFold}(w_1, x))$ и тождественно сохраняет наблюдаемую $\pi(\mathrm{stageFold}(w, x)) = \pi(x)$, а для любой проективной системы гомоморфизмов $\pi_n : M_{n+1} \to M_n$ ($\mathrm{obs}_{n+1}(x) = \mathrm{obs}_n(\pi_n(x))$) композиция проекций на базовый носитель $\Pi_{n\to 0} : M_n \to M_0$ стабилизирует все макроскопические наблюдаемые на любом этаже~$n\in\mathbb{N}$:
  \begin{equation}
    \label{eq:mono-tower-stabilization}
    \mathrm{obs}_0\bigl(\Pi_{n\to 0}(x_n)\bigr) \;=\; \mathrm{obs}_n(x_n) \qquad \forall\,n\in\mathbb{N},\ \forall\,x_n\in M_n.
  \end{equation}
\item \textbf{Когомологическая аннигиляция ядра $\ker\pi \subseteq \ker d_1 \cap \ker d_2$ и телескопирование одномерной границы (\texttt{boundarySum}):}
  всякое калибровочное возмущение из ядра $\ker\pi$, раскладывающееся в узловую кограницу $d_0(\delta\chi)$ на рёбрах и рёберную кограницу $d_1(\delta q)$ на гранях, тождественно аннигилирует во всех физических наблюдаемых сорта~$\mathcal{T}$, а сумма первых разностей вдоль любой конечной цепи стадий длины $N\in\mathbb{N}$ телескопируется к граничным точкам:
  \begin{align}
    d_1\bigl(r + d_0(\delta\chi)\bigr) &\equiv d_1(r), \qquad
    d_2\bigl(n + d_1(\delta q)\bigr) \;\equiv\; d_2(n), \label{eq:mono-kernel-annihilation} \\
    \sum_{k=0}^{N-1}\bigl(F(k+1) - F(k)\bigr) &= F(N) - F(0). \label{eq:mono-boundary-telescoping}
  \end{align}
\item \textbf{Конструктивная теорема о сжатии двусторонней вилки (\texttt{bracket\allowbreak FoldL}, \texttt{bracket\allowbreak FoldU}):}
  для любого списка неотрицательных сужений стадий $w = [\delta_1, \dots, \delta_r]$ ($\delta_i \ge 0$) монотонно возрастающая нижняя граница $L_w = L_0 + \sum_i \delta_i$ и монотонно убывающая верхняя граница $U_w = U_0 - \sum_i \delta_i$, охватывающие предел $L_w \le X \le U_w$, при обращении ширины вилки в нуль однозначно фиксируют предел по антисимметрии целочисленного порядка в~$\mathbb{Z}$:
  \begin{equation}
    \label{eq:mono-bracket-squeeze}
    L_w \le X \le U_w \;\land\; U_w - L_w = 0 \quad\Longrightarrow\quad L_w \;=\; X \;=\; U_w.
  \end{equation}
\item \textbf{Двусторонняя стереографическая эквивалентность бесконечного счётного носителя $\mathcal{M}$ и континуального носителя $\mathcal{T}$:}
  дискретный носитель~$\mathcal{M}$ есть бесконечное счётное индуктивное семейство колец $\mathbb{Z}_{4Q}$ по всем масштабам четверть-периода $Q = C + S$, всем виткам обмотки $q\in\mathbb{Z}$ и узлам $\mathbb{Z}^d$ (где на масштабе $K=1$ фактор-регистр одного оборота одной элементарной ячейки равен $\mathbb{Z}_{512}$). Рациональный стереографический подъём Вейерштрасса $\iota_{T\to M}(C,S) = (C^2-S^2,\ 2CS,\ C^2+S^2)$ и обратное проектирование $\pi_{M\to T}$ связывают $\mathcal{M}$ и $\mathcal{T}$ точным пифагоровым тождеством $(C^2-S^2)^2 + (2CS)^2 \equiv (C^2+S^2)^2$, постоянным разностным вронскианом $C\bigl((S+1)-S\bigr) - S\bigl((C-1)-C\bigr) \equiv C + S = Q$ и взаимно-обратными соотношениями $(C+S) + (C-S) = 2C$, $(C+S) - (C-S) = 2S$ на всех масштабах~$Q$.
\end{enumerate}

\section{Конструктивная непротиворечивость и теорема Ловера--Гёделя--Тарского}

\begin{theorem}[Общая конструкция предела, конструктивная непротиворечивость и необходимость внешнего экспериментального арбитра, модуль \texttt{ModelTheoreticLimit.lean}]
\label{thm:mono-godel-tarski}
В двусортной системе носителей $\langle\mathcal{M},\mathcal{T};\,\pi_{M\to T},\,\iota_{T\to M}\rangle$:
\begin{enumerate}[label=\textbf{(\arabic*)},leftmargin=*]
\item \textbf{Отсутствие парадоксов неконструктивного выбора и разрешение булевой самореференции во времени:}
  поскольку каждое локальное звено счётного индуктивного носителя~$\mathcal{M}$ (и его однооборотный фактор-регистр $\mathbb{Z}_{512}$ на масштабе $K=1$) снабжено разрешимым равенством (\texttt{DecidableEq}, \texttt{0 sorryAx}, \texttt{0 Classical.choice}), на нём невозможны неконструктивные разбиения Банаха--Тарского, биекция одноразового блокнота Вернама--Шеннона $k = (c - m)\bmod 512 \in \mathbb{Z}_{512}$ сохраняет $H(\mathcal{M}\mid\mathcal{T}) = H(\mathcal{M})$, а статический парадокс булевой самонегации ($\forall b\in\mathrm{Bool},\ b \ne \neg b$) разрешается как полуволновой осциллятор периода~$2$: $\varphi \mapsto (\varphi + 256)\bmod 512$ без неподвижных точек.
\item \textbf{Полиморфная диагональная теорема Ловера--Гёделя--Тарского и теоретико-модельная невыразимость семантического оракула истинности формул:}
  для любого типа синтаксических кодов $\alpha : \mathrm{Type}$, любого типа семантических истинностных значений $\beta : \mathrm{Type}$, любого внутреннего вычислителя $\mathrm{eval} : \alpha \to (\alpha \to \beta)$ и любой семантической инверсии без неподвижных точек $\sigma : \beta \to \beta$ ($\forall b:\beta,\ \sigma(b) \ne b$), внешний семантический оракул $\mathrm{oracle}(x) := \sigma(\mathrm{eval}\ x\ x)$ строго невыразим никаким внутренним синтаксическим кодом $p : \alpha$, а для любого языка формул $\mathrm{Formula}$ с отрицанием $\mathrm{neg}$, подстановкой кодов $\mathrm{sub} : \mathrm{Code}\to\mathrm{Code}\to\mathrm{Formula}$ и внешним семантическим отношением истинности $\mathrm{oracleModels} : \mathrm{Formula}\to\mathrm{Prop}$ ($\forall\varphi,\ \mathrm{oracleModels}(\mathrm{neg}\,\varphi)\leftrightarrow\neg\mathrm{oracleModels}(\varphi)$) ни один внутренний предикат $\mathrm{internalEval}(c,\cdot)$, кодирующий диагональное отрицание, не может совпадать с внешним оракулом на $\mathrm{sub}(c,c)$:
  \begin{align}
    \mathrm{eval}(p, p) &\ne \mathrm{oracle}(p) \qquad (\text{для любого } p : \alpha), \label{eq:mono-lawvere-godel-tarski} \\
    \mathrm{internalEval}\bigl(c,\, \mathrm{sub}(c, c)\bigr) &\not\Longleftrightarrow \mathrm{oracleModels}\bigl(\mathrm{sub}(c, c)\bigr). \label{eq:mono-tarski-undefinability}
  \end{align}
\end{enumerate}
\end{theorem}

\begin{proof}
В пункте~(1) при $\pi(x) = \pi(x')$ и $\pi(y) = \pi(y')$ из гомоморфности $\pi(x \star_A y) = \pi(x)\star_B\pi(y) = \pi(x')\star_B\pi(y') = \pi(x' \star_A y')$ следует конгруэнтность ядра~\eqref{eq:mono-malcev-congruence}, по индукции по списку стадий $w$ и номеру этажа $n$ следует стабилизация проективной башни~\eqref{eq:mono-tower-stabilization}, а из $d_1\circ d_0 \equiv 0$ и $d_2\circ d_1 \equiv 0$ следует аннигиляция ядра~\eqref{eq:mono-kernel-annihilation}. На $\mathbb{Z}_{512}$ сдвиг $\varphi \mapsto (\varphi + 256)\bmod 512$ удовлетворяет $(\varphi + 256) \not\equiv \varphi \pmod{512}$ и $(\varphi + 256 + 256) \equiv \varphi \pmod{512}$.

В пункте~(2), если бы существовал код $p : \alpha$, для которого $\mathrm{eval}\ p = \mathrm{oracle}$, то подстановка диагонального аргумента $x = p$ давала бы $\mathrm{eval}\ p\ p = \mathrm{oracle}\ p = \sigma(\mathrm{eval}\ p\ p)$, что противоречит отсутствию неподвижных точек $\forall b:\beta,\ \sigma(b) \ne b$. Для формулы $\psi := \mathrm{sub}(c,c)$ условие диагонального кодирования $\mathrm{internalEval}(c,\psi) \Longleftrightarrow \neg\mathrm{oracleModels}(\psi)$ вместе с предположением $\mathrm{internalEval}(c,\psi) \Longleftrightarrow \mathrm{oracleModels}(\psi)$ даёт $\mathrm{oracleModels}(\psi) \Longleftrightarrow \neg\mathrm{oracleModels}(\psi)$, что конструктивно невозможно.
\end{proof}

Совместное действие теоремы Мальцева о ядре гомоморфизма~\eqref{eq:mono-malcev-congruence}--\eqref{eq:mono-kernel-annihilation} и диагональной теоремы Ловера--Гёделя--Тарского~\eqref{eq:mono-lawvere-godel-tarski} даёт строгий математический ответ на традиционный скепсис по отношению к «теориям всего»:
\begin{enumerate}[label=\textbf{\arabic*.},leftmargin=*]
\item \textbf{Почему невозможно «следующее приближение» $M_{*+1}$ (теорема Мальцева об обрыве проективной башни на самозамкнутом регистре ячейки):}
  пока физические теории строились над континуальным фоном, точка которого предполагала бесконечный запас бит, любая модель $M_n$ оставалась лишь промежуточным приближением, допускающим бесконечное достраивание новых уровней $M_{n+1}$ и новых подгоночных параметров. Однако неделимость элементарного объёма ($1$ состояние собственной занятости центра) и плотная упаковка всех ячеек без пустот ($1 + 2^d = d^2 \iff d = 3$) однозначно фиксируют 9-битный однооборотный регистр ячейки $\mathbb{Z}_{512}$ и оболочку кубооктаэдра $O_2(3)$, а двустороннее проективно-ретракционное соответствие Вейерштрасса $\pi_{M\to T} \leftrightarrow \iota_{T\to M}$ замыкает носители $\mathcal{M} \simeq \mathcal{T}$. По формулам~\eqref{eq:mono-tower-stabilization}--\eqref{eq:mono-kernel-annihilation} для любого гипотетического «более тонкого» уширения $\pi : M_{*+1} \twoheadrightarrow \mathcal{M}$ всё его ядро $\ker\pi \subseteq \ker d_1 \cap \ker d_2$ тождественно аннигилирует во всех физических наблюдаемых. Следовательно, цепочка физических приближений математически обрывается на $\lim_{n\to\infty} M_n = \mathcal{M}$: никакого «следующего приближения» за пределами этого носителя не существует.
\item \textbf{Почему единственным судьёй теории является прямой лабораторный эксперимент (рубеж Гёделя--Тарского при нуле свободных параметров):}
  традиционные умозрительные «теории всего» страдают либо нефальсифицируемостью из-за десятков свободных параметров и ландшафтов из $10^{500}$ вакуумов (позволяющих подогнать формулы под любой исход опыта), либо претензией на чисто априорную истинность без эксперимента. В построенной здесь системе оба этих пути закрыты доказательно:
  \begin{itemize}[noitemsep]
  \item во всей дедуктивной цепи от $1 + 2^d = d^2$ до $\alpha_{\mathrm{fs}}^{-1}(3) \in \mathbb{Q}(\sqrt{2})$ и порога фотокаскада $120 + 8 = 128$ \textbf{ровно ноль свободных параметров}, а внутренняя дедуктивная непротиворечивость полностью удостоверена ядром Lean~4 (\texttt{0 sorryAx}, \texttt{0 Classical.choice});
  \item в то же время по теореме Ловера--Гёделя--Тарского~\eqref{eq:mono-lawvere-godel-tarski} никакая непротиворечивая формальная система не может изнутри собственного синтаксиса доказать своё семантическое тождество с внешней физической природой, а поскольку в точке предела $\lim_{n\to\infty} M_n = \mathcal{M}$ переход к следующей синтаксической надстройке $M_{*+1}$ невозможен, \textbf{единственным семантическим оракулом (судьёй) остаётся прямой фальсифицируемый лабораторный эксперимент}, перед которым у теории без единого подгоночного параметра нет люфта для уклонения.
  \end{itemize}
\end{enumerate}

Именно это эпистемологическое разделение труда между математическим доказательством непротиворечивости и лабораторным оракулом определяет четырёхтомную архитектуру настоящей серии:
\begin{itemize}[leftmargin=*]
\item \textbf{Теория (настоящая Книга~1, \texttt{formal-foundations-monograph.pdf}):} замкнутый дедуктивный вывод законов природы и констант на $\langle\mathcal{M},\mathcal{T};\,\pi_{M\to T},\,\iota_{T\to M}\rangle$ с нулём свободных параметров и машинной проверкой непротиворечивости в Lean~4;
\item \textbf{Лаборатория (Книга~2, \texttt{dqm-laboratory-manual.pdf}):} внешний семантический арбитраж по теореме Гёделя--Тарского --- пошаговый протокол чистой комнаты для четырёхслойного бесфононного кремниевого гетероперехода «Кварц -- Тетрацен -- $\mathrm{SiO_x}$ -- IBC-$\mathrm{Si}$» ($Z(\mathrm{C_{18}H_{12}}) + Z(\mathrm{O}) = 120 + 8 = 128 = Q_\square(3)$, барьер пассивации Паули $45^2 < 2\,048 < 46^2$), контрольный A/B-тест порога Паули ($Z=45$ против $Z=46$), четыре сигнатуры фальсификации (\texttt{M1}--\texttt{M4}), упреждающая матрица коррекции шести технологических режимов срыва (\texttt{F1}--\texttt{F6}) и журнал критических экспериментов;
\item \textbf{Инженерия (Книга~3, \texttt{dqm-engineering-handbook.pdf}):} теоретическая электротехника фазовой эстафеты ($V = \mathcal{C}_3(\Delta\varphi),\ I = \mathcal{S}_3(\Delta\varphi)$), беспульсационные квадратурные выпрямители, физическая реализация вычислений без тепла Ландауэра в коридоре Гейзенберга $|\Delta\varphi|\le 40 < 41$, безветвевой кремниевый стек \texttt{Ring9} / \texttt{Bus36} / \texttt{VerletCell72} и 12-катушечные резонаторы кубооктаэдра $O_2(3)$;
\item \textbf{Научно-популярный путеводитель (Книга~4, \texttt{natural-sciences-course.pdf}, \emph{Powered by Dr.\,STONE}):} последовательный разговорный курс всех семи ступеней крафта Вселенной.
\end{itemize}
